% !TeX root = ../despliegue.tex
% ===========================================================================
%  Presentación del documento. Va antes del índice, sin numerar.
% ===========================================================================

\section*{Sobre este manual}
\addcontentsline{toc}{section}{Sobre este manual}
\markboth{Sobre este manual}{}

Este documento corresponde al \textbf{entregable E5} del proyecto integrador y
describe cómo levantar, verificar y operar el Sistema de Reserva de Canchas
Deportivas.

Se ha escrito pensando en quien recibe el proyecto sin haber participado en él:
alguien que clona el repositorio, quiere verlo funcionando en su propia máquina
y necesita saber qué hacer cuando algo no arranca. Por eso el texto va directo a
los comandos y deja fuera las justificaciones de diseño.

Los otros dos documentos del proyecto cubren lo que aquí se omite:

\begin{itemize}
  \item El \textbf{informe} (entregable E1) explica la arquitectura, el modelo
        de datos y las reglas de negocio: \emph{qué} hace cada pieza y
        \emph{por qué} está construida así.
  \item El \textbf{manual de usuario} explica cómo se opera la aplicación desde
        el navegador, tanto para el usuario final como para el administrador.
\end{itemize}

\nota{Todo lo que sigue está verificado sobre la versión del sistema entregada.
Los comandos se ejecutan desde la raíz del repositorio, salvo que se indique
otra cosa.}
